\section{Traditional Metaphysics}

\begin{quotex}
Guenon's \emph{Man and his Becoming according to the Vedanta} will draw the attention of the well trained and qualified reader. Of course, it will also become the source of misunderstandings for a certain category of ``third rate'' critics and intellectuals who oscillate between platitudes and political and spiritual fancies.
\flright{\textsc{Julius Evola}}
\end{quotex}


Rene Guenon was not happy with Nicholson's translation of \emph{Man and his Becoming}. A literal word for word translation of a text that may read well in French, may not result in smooth English. Many people have informed me that they have had trouble understanding \emph{Man and his Becoming}. That may be due, in part, not to the difficulty of the concepts but rather to the confusion of the translation. Hence, I have begun a slow-motion translation of the most important sections of his major works of metaphysics.

This is the beginning of a slow motion translation, still in progress, of the fundamental texts of metaphysics as described by Rene Guenon. These texts are mostly found in the following books:

\begin{itemize}


\item \emph{Man and his Becoming according to the Vedanta}


\item \emph{The Symbolism of the Cross}


\item \emph{The Multiple States of the Being}

\end{itemize}


The first goal of the translations is to make it read smoothly in English; this is not possible with a word for word translation, which often is misleading. In this way, the fundamental points of the text.

The second goal is to employ the most commonly accepted English words to express the metaphysical ideas. This is often not the literal translation of the French word.

The third goal is to make the style and terminology consistent among the three works listed above.

The first text selected is \emph{Chapter II: Fundamental Distinction between Self and Ego} from \emph{Man and his Becoming according to the Vedanta}. I will be updating it depending on my time commitments and positive feedback. Here is a teaser with a fresh translation of Chapter 2 on the Self and the ego.

See Principles Of Metaphysics\footnote{\url{https://www.gornahoor.net/library/PrinciplesOfMetaphysics.pdf}}


\flrightit{Posted on 2022-09-03 by Cologero}

